% Captions only; use the corresponding PDF with \includegraphics.
% Do not imply analytical certification or successful mismatch recovery.
\newcommand{\ExperimentOverviewCaption}{Information timing changes collective adaptation in a finite-action stochastic aggregative game. (A) Action-value gaps under common and staggered delivery; positive values favor activation. (B) Corresponding aggregate activation. Panels A--B use the same median-selected illustrative seed. (C) Mean excess buildup across all tested delivery rules; common randomized timing can achieve the demonstrated reduction without agent-specific policy distributions. (D) A training-selected correction based on aggregate feedback does not reliably recover performance under memory mismatch. Panel C uses 96 held-out seeds; D uses 12 independent batches of four probe and four fresh evaluation episodes per memory. Intervals are pointwise 95\% bootstrap intervals. Dashed lines mark an empirical mean target, not a certificate.}
\newcommand{\ExperimentMechanismCaption}{Beliefs, action-value gaps, binary activation and physical response for matched seed 1003. Faint belief traces are individual agents; heavy curves are population means. Heatmap values are clipped to $[-0.6,0.6]$ for display; white completed-service cells have no active choice. Vertical lines mark full-information follow-up beginning at round 31. The seed minimizes distance to the median paired excess-buildup scores, rather than maximizing the intervention effect.}
\newcommand{\ExperimentDesignCaption}{Policy design and decision quality on 96 held-out seeds. (A) Every tested fixed-frequency policy; circles pass the empirical quality/service screen and crosses fail it. (B) Marginal changes from adding group-phase interventions, colored by group, showing that further changes can hurt. (C) Worst-agent average complete-history reference decision loss. Low excess buildup alone does not establish good decisions. Reference loss is neither equilibrium regret nor actual welfare loss.}
\newcommand{\ExperimentFeedbackCaption}{Aggregate-feedback correction under memory mismatch. Four independent reset episodes form each probe; the selected policy is evaluated on four fresh episodes, with 12 batches per memory. (A) Held-out excess buildup. (B) Compatibility-set outcomes with 95\% Wilson intervals over independent decisions. (C) Differentiation of selected policies. Known-memory selection uses the true label and training data, not hindsight optimization. Empty compatibility sets use the conservative fallback.}
\newcommand{\ExperimentControlsCaption}{Paired policy-minus-common-periodic differences across six conditions, with 64 independent seeds per condition and pointwise 95\% bootstrap intervals. Negative differences favor the intervention. Zero pressure cost removes only the state-dependent cost term. The omitted-public-inference ablation changes the learner's information processing; it is not the primary observation model.}
